% Sorgente LaTeX canonico, importato e revisionato dal precedente sorgente Markdown.
\chapter{Successioni numeriche}
\label{chap:6}

\begin{obiettivocapitolo}{calcolare limiti e comportamento di successioni esplicite o ricorrenti}
\prioritaalta. Prima di applicare un criterio, identifica se la successione è esplicita, monotona, ricorrente, oscillante o dominata da fattoriali/esponenziali.
\end{obiettivocapitolo}
\begin{sceltametodo}{limite di una successione}
\begin{tabularx}{\textwidth}{@{}p{.30\textwidth}X@{}}
Rapporto di polinomi in $n$ & dividi per la massima potenza di $n$ o confronta i gradi.\\
Radici con differenze & razionalizza oppure raccogli il termine dominante.\\
Potenze/esponenziali/fattoriali & usa gerarchie di crescita, rapporto o logaritmo.\\
Forma $(1+u_n)^{v_n}$ & passa al logaritmo e studia $v_n\log(1+u_n)$.\\
Ricorrenza $a_{n+1}=\varphi(a_n)$ & prova monotonia e limitatezza; poi il limite soddisfa $L=\varphi(L)$.\\
Oscillazione & cerca sottosuccessioni con limiti diversi o calcola $\liminf$ e $\limsup$.\\
Differenza/rapporto adatto & considera Stolz--Cesàro solo dopo averne verificato le ipotesi.
\end{tabularx}
\end{sceltametodo}
\begin{proceduraesame}{successione definita per ricorrenza}
\begin{enumerate}
\item Dimostra che tutti i termini sono definiti e restano in un intervallo invariante.
\item Studia il segno di $a_{n+1}-a_n$ oppure confronta $\varphi(x)$ con $x$ per ottenere monotonia.
\item Trova un maggiorante o minorante indipendente da $n$.
\item Concludi la convergenza tramite monotonia e limitatezza.
\item Risolvi $L=\varphi(L)$ e scegli il punto fisso compatibile con l'intervallo dei termini.
\end{enumerate}
\end{proceduraesame}
\begin{errorefrequente}{punto fisso non dimostra convergenza}
Risolvere $L=\varphi(L)$ fornisce soltanto i possibili limiti. Senza una prova di convergenza, la successione può oscillare, divergere o non essere definita per ogni indice.
\end{errorefrequente}
\begin{controllofinale}{successioni}
Confronta il risultato con segno, monotonia, limiti inferiori/superiori e primi termini numerici.
\end{controllofinale}

Una successione \((a_n)\) è una funzione il cui argomento è un indice naturale.

\section{Definizione e indicizzazione}\label{definizione-e-indicizzazione}

\[
a:\mathbb N_0\to\mathbb R,
\qquad
n\mapsto a_n.
\]

\[
\exists N_0\in\mathbb N_0:\
\forall n\ge N_0,\ a_n=b_n
\Longrightarrow
\left(a_n\to L\Longleftrightarrow b_n\to L\right).
\]

\[
a_n\to L
\Longleftrightarrow
a_{n+n_0}\to L,
\qquad n_0\in\mathbb N_0.
\]

Approfondimento: \href{https://ingegnerismo.it/matematica/successione/}{successione}.

\section{Successioni elementari}\label{successioni-elementari}

Per \(q\in\mathbb R\):

\[
\lim_{n\to\infty}q^n=
\begin{cases}
0,&|q|<1,\\
1,&q=1,\\
+\infty,&q>1,
\end{cases}
\]

mentre per \(q\le-1\) la successione non converge; se \(q<-1\) il modulo tende a \(+\infty\) e il segno alterna.

Per \(\alpha\in\mathbb R\), considerando \(n\ge1\):

\[
\lim_{n\to\infty}n^\alpha=
\begin{cases}
+\infty,&\alpha>0,\\
1,&\alpha=0,\\
0,&\alpha<0.
\end{cases}
\]

\section{Limite finito}\label{limite-finito}

\[
\lim_{n\to\infty}a_n=L
\Longleftrightarrow
\forall\varepsilon>0,\ \exists N\in\mathbb N_0:\
n\ge N\Rightarrow |a_n-L|<\varepsilon.
\]

Forma equivalente:

\[
\forall\varepsilon>0,\ \exists N:\
n\ge N\Rightarrow L-\varepsilon<a_n<L+\varepsilon.
\]

Negazione:

\[
a_n\not\to L
\Longleftrightarrow
\exists\varepsilon_0>0,\ \forall N,\ \exists n\ge N:\
|a_n-L|\ge\varepsilon_0.
\]

Approfondimenti: \href{https://ingegnerismo.it/matematica/limite-di-successione/}{limite di successione} e \href{https://ingegnerismo.it/matematica/unicita-del-limite/}{unicità del limite}.

\section{Unicità del limite}\label{unicituxe0-del-limite}

\[
a_n\to L,
\qquad
a_n\to L'
\Longrightarrow
L=L'.
\]

Approfondimento: \href{https://ingegnerismo.it/matematica/unicita-del-limite/}{unicità del limite}.

\section{Limitatezza delle successioni convergenti}\label{limitatezza-delle-successioni-convergenti}

\[
a_n\to L\in\mathbb R
\Longrightarrow
\exists M>0:\ |a_n|\le M\quad\forall n.
\]

\[
a_n\text{ limitata}
\Longleftrightarrow
\exists m,M\in\mathbb R:\ m\le a_n\le M\quad\forall n.
\]

Il converso non vale.

Approfondimento: \href{https://ingegnerismo.it/matematica/limitatezza-delle-successioni-convergenti/}{limitatezza delle successioni convergenti}, \href{https://ingegnerismo.it/matematica/successione/}{successione}.

\section{Permanenza del segno}\label{permanenza-del-segno}

\[
a_n\to L>0
\Longrightarrow
a_n>\frac L2>0
\quad\text{definitivamente},
\]

\[
a_n\to L<0
\Longrightarrow
a_n<\frac L2<0
\quad\text{definitivamente}.
\]

\[
a_n\ge0\quad\text{definitivamente},
\qquad
a_n\to L
\Longrightarrow
L\ge0.
\]

\[
a_n\le0\quad\text{definitivamente},
\qquad
a_n\to L
\Longrightarrow
L\le0.
\]

Approfondimento: \href{https://ingegnerismo.it/matematica/permanenza-del-segno-per-successioni-convergenti/}{permanenza del segno per successioni convergenti}.

\section{Limite infinito}\label{limite-infinito}

\[
a_n\to+\infty
\Longleftrightarrow
\forall M>0,\ \exists N:\ n\ge N\Rightarrow a_n>M,
\]

\[
a_n\to-\infty
\Longleftrightarrow
\forall M>0,\ \exists N:\ n\ge N\Rightarrow a_n<-M.
\]

Negazioni:

\[
a_n\not\to+\infty
\Longleftrightarrow
\exists M_0>0,\ \forall N,\ \exists n\ge N:\ a_n\le M_0,
\]

\[
a_n\not\to-\infty
\Longleftrightarrow
\exists M_0>0,\ \forall N,\ \exists n\ge N:\ a_n\ge-M_0.
\]

Approfondimento: \href{https://ingegnerismo.it/matematica/limite-di-successione/}{limite di successione}.

\section{Algebra dei limiti}\label{algebra-dei-limiti}

Se \(a_n\to a\), \(b_n\to b\) con \(a,b\in\mathbb R\) e \(\lambda\in\mathbb R\):

\[
a_n\pm b_n\to a\pm b,
\qquad
\lambda a_n\to\lambda a,
\]

\[
a_nb_n\to ab,
\qquad
|a_n|\to|a|,
\]

\[
\frac{a_n}{b_n}\to\frac ab,
\qquad
b\ne0,
\]

poiché \(b_n\ne0\) definitivamente per la permanenza del segno.

\[
a_n^k\to a^k,
\qquad k\in\mathbb N_+.
\]

\[
a_n\to a\ne0,
\qquad
k\in\mathbb Z_{<0}
\Longrightarrow
a_n^k\to a^k.
\]

\[
a_n\to a,
\qquad
a_n\ge0\text{ definitivamente},
\qquad
a\ge0
\Longrightarrow
\sqrt[2m]{a_n}\to\sqrt[2m]{a},
\qquad m\in\mathbb N_+.
\]

\[
a_n\to a
\Longrightarrow
\sqrt[2m+1]{a_n}\to\sqrt[2m+1]{a},
\qquad m\in\mathbb N_0.
\]

Approfondimento: \href{https://ingegnerismo.it/matematica/algebra-dei-limiti-per-successioni-convergenti/}{algebra dei limiti per successioni convergenti}.

\section{Ordine, confronto e teorema dei carabinieri}\label{ordine-confronto-e-teorema-dei-carabinieri}

Nelle prime due formule \(A,B\in\mathbb R\).

\[
a_n\le b_n\quad\text{definitivamente},
\qquad
a_n\to A,
\qquad
b_n\to B
\Longrightarrow
A\le B.
\]

\[
a_n\to A,
\qquad
b_n\to B,
\qquad
A<B
\Longrightarrow
a_n<b_n\quad\text{definitivamente}.
\]

\[
\ell_n\le x_n\le u_n\quad\text{definitivamente},
\qquad
\ell_n,u_n\to L
\Longrightarrow
x_n\to L.
\]

\[
|x_n-L|\le r_n,
\qquad
r_n\to0
\Longrightarrow
x_n\to L.
\]

\[
a_n\le b_n\quad\text{definitivamente},
\qquad
a_n\to+\infty
\Longrightarrow
b_n\to+\infty,
\]

\[
a_n\le b_n\quad\text{definitivamente},
\qquad
b_n\to-\infty
\Longrightarrow
a_n\to-\infty.
\]

Approfondimento: \href{https://ingegnerismo.it/matematica/ordine-e-confronto-per-limiti-di-successioni/}{ordine e confronto per limiti di successioni}, \href{https://ingegnerismo.it/matematica/teorema-dei-carabinieri/}{teorema dei carabinieri}.

\section{Successioni monotone}\label{successioni-monotone}

\[
a_{n+1}\ge a_n\quad\forall n
\Longleftrightarrow
(a_n)\text{ non decrescente},
\]

\[
a_{n+1}\le a_n\quad\forall n
\Longleftrightarrow
(a_n)\text{ non crescente}.
\]

\[
\begin{array}{c|c|c}
\text{monotonia}&\text{vincolo}&\text{limite}\\
\hline
\nearrow&\text{superiormente limitata}&\sup_n a_n\\
\nearrow&\text{non superiormente limitata}&+\infty\\
\searrow&\text{inferiormente limitata}&\inf_n a_n\\
\searrow&\text{non inferiormente limitata}&-\infty
\end{array}
\]

Approfondimento: \href{https://ingegnerismo.it/matematica/successione-monotona/}{successione monotona}, \href{https://ingegnerismo.it/matematica/completezza-dei-reali/}{completezza dei reali}.

\section{Successioni definite per ricorrenza}\label{successioni-definite-per-ricorrenza}

\[
a_{n+1}=\varphi(a_n),
\qquad a_0\text{ assegnato}.
\]

\[
a_n\to L,
\qquad
\varphi\text{ continua in }L
\Longrightarrow
L=\varphi(L).
\]

Se \(I\) è chiuso, \(a_0\in I\), \(\varphi(I)\subseteq I\) e \(\varphi\) è una contrazione, cioè

\[
|\varphi(x)-\varphi(y)|\le q|x-y|,
\qquad
0\le q<1,
\qquad x,y\in I,
\]

allora esiste un unico punto fisso \(L\in I\) e

\[
a_n\to L,
\]

\[
|a_n-L|\le q^n|a_0-L|,
\]

\[
|a_n-L|\le\frac{q^n}{1-q}|a_1-a_0|.
\]

Ricorrenza affine:

\[
a_{n+1}=qa_n+b.
\]

Per \(q\ne1\):

\[
L=\frac b{1-q},
\qquad
a_n=L+q^n(a_0-L).
\]

\[
\begin{array}{c|c}
q&\text{comportamento per }q\ne1,\ a_0\ne L\\
\hline
|q|<1&a_n\to L\\
q=-1&\text{alternanza di due valori}\\
q>1&|a_n|\to\infty\text{ con segno determinato}\\
q<-1&\text{oscillazione con modulo crescente}
\end{array}
\]

\[
q\ne1,\quad a_0=L
\Longrightarrow
a_n=L.
\]

\[
q=1
\Longrightarrow
a_n=a_0+nb,
\qquad
\lim_{n\to\infty}a_n=
\begin{cases}
a_0,&b=0,\\
+\infty,&b>0,\\
-\infty,&b<0.
\end{cases}
\]

Approfondimento: \href{https://ingegnerismo.it/matematica/successioni-ricorsive/}{successioni ricorsive}.

\section{Numero di Nepero}\label{numero-di-nepero}

\[
e=\lim_{n\to\infty}\left(1+\frac1n\right)^n
=\sum_{k=0}^{\infty}\frac1{k!}
\approx2{,}718281828.
\]

\[
\left(1+\frac xn\right)^n\to e^x,
\qquad x\in\mathbb R.
\]

\[
\left(1+\frac1n\right)^n<e<\left(1+\frac1n\right)^{n+1},
\qquad n\ge1.
\]

Approfondimento: \href{https://ingegnerismo.it/matematica/numero-di-nepero/}{numero di Nepero}.

\section{Sottosuccessioni e Bolzano--Weierstrass}\label{sottosuccessioni-e-bolzanoweierstrass}

\[
a_{n_k},
\qquad
n_0<n_1<\cdots<n_k<\cdots.
\]

\[
a_n\to L
\Longrightarrow
a_{n_k}\to L
\quad\text{per ogni sottosuccessione}.
\]

\[
\text{ogni successione reale ammette una sottosuccessione monotona}.
\]

\[
(a_n)\text{ limitata}
\Longrightarrow
\exists(a_{n_k})\text{ convergente}.
\]

\[
a_{n_k}\to L_1,
\qquad
a_{m_k}\to L_2,
\qquad
L_1\ne L_2
\Longrightarrow
(a_n)\text{ non converge}.
\]

Approfondimenti: \href{https://ingegnerismo.it/matematica/sottosuccessione/\#teorema-della-sottosuccessione-monotona}{teorema della sottosuccessione monotona} e \href{https://ingegnerismo.it/matematica/teorema-di-bolzano-weierstrass/}{teorema di Bolzano--Weierstrass}.

Esercizi svolti: \href{https://ingegnerismo.it/matematica/sottosuccessioni-bolzano-esercizi/}{sottosuccessioni e Bolzano--Weierstrass}.

\section{Limite superiore e limite inferiore}\label{limite-superiore-e-limite-inferiore}

Le quantità seguenti sono considerate nella retta reale estesa \(\overline{\mathbb R}\):

\[
U_N=\sup_{n\ge N}a_n,
\qquad
L_N=\inf_{n\ge N}a_n.
\]

\[
U_{N+1}\le U_N,
\qquad
L_{N+1}\ge L_N.
\]

\[
\limsup_{n\to\infty}a_n
=\lim_{N\to\infty}U_N
=\inf_N U_N,
\]

\[
\limsup_{n\to\infty}a_n
=\inf_{N}\sup_{n\ge N}a_n,
\]

\[
\liminf_{n\to\infty}a_n
=\lim_{N\to\infty}L_N
=\sup_N L_N,
\]

\[
\liminf_{n\to\infty}a_n
=\sup_{N}\inf_{n\ge N}a_n.
\]

\[
\liminf a_n\le\limsup a_n.
\]

Quando almeno uno dei due estremi non è finito, i soli profili compatibili con \(\liminf a_n\le\limsup a_n\) sono i cinque seguenti (con \(L\in\mathbb R\)):

\begin{longtable}[]{@{}rrl@{}}
\toprule\noalign{}
\(\liminf a_n\) & \(\limsup a_n\) & Esempio \\
\midrule\noalign{}
\endhead
\bottomrule\noalign{}
\endlastfoot
\(-\infty\) & \(-\infty\) & \(a_n=-n\) \\
\(-\infty\) & \(L\) & \(a_{2k}=L\), \(a_{2k+1}=-(k+3)\) \\
\(-\infty\) & \(+\infty\) & \(a_n=(-1)^n n\) \\
\(L\) & \(+\infty\) & \(a_{2k}=L\), \(a_{2k+1}=k+3\) \\
\(+\infty\) & \(+\infty\) & \(a_n=n\) \\
\end{longtable}

\[
a_n\to L\in\overline{\mathbb R}
\Longleftrightarrow
\liminf a_n=\limsup a_n=L.
\]

Per \((a_n)\) limitata:

\[
\liminf a_n=\min\{\text{valori di aderenza}\},
\qquad
\limsup a_n=\max\{\text{valori di aderenza}\}.
\]

Approfondimento: \href{https://ingegnerismo.it/matematica/limite-superiore-e-limite-inferiore-di-una-successione/}{limite superiore e limite inferiore di una successione}.

Esercizi svolti: \href{https://ingegnerismo.it/matematica/limsup-liminf-successioni-esercizi/}{limite superiore e limite inferiore}.

\section{Criterio di Cauchy}\label{criterio-di-cauchy}

\[
(a_n)\text{ di Cauchy}
\Longleftrightarrow
\forall\varepsilon>0,\ \exists N\in\mathbb{N},\ \forall m,n\in\mathbb N:\
m,n\ge N\Longrightarrow |a_n-a_m|<\varepsilon.
\]

Negazione:

\[
\exists\varepsilon_0>0,\ \forall N\in\mathbb N,\ \exists m,n\in\mathbb N:\
m,n\ge N\ \land\ |a_n-a_m|\ge\varepsilon_0.
\]

In \(\mathbb R\):

\[
a_n\text{ converge}
\Longleftrightarrow
(a_n)\text{ è di Cauchy}.
\]

\[
(a_n)\text{ di Cauchy}
\Longrightarrow
(a_n)\text{ limitata}.
\]

Approfondimento: \href{https://ingegnerismo.it/matematica/criterio-di-cauchy/}{criterio di Cauchy}, \href{https://ingegnerismo.it/matematica/completezza-dei-reali/}{completezza dei reali}.

\section{Limiti notevoli e gerarchie}\label{limiti-notevoli-e-gerarchie}

\[
\left(1+\frac xn\right)^n\to e^x,
\qquad x\in\mathbb R,
\]

\[
\sqrt[n]a\to1,
\qquad a>0,
\]

\[
\sqrt[n]n\to1,
\]

\[
\frac{n^\alpha}{a^n}\to0,
\qquad a>1,
\qquad \alpha>0,
\]

\[
\frac{a^n}{n!}\to0,
\qquad a>0.
\]

\[
\ln n\ll n^\alpha\ll a^n\ll n!\ll n^n,
\qquad
\alpha>0,
\qquad
a>1.
\]

Formula di Stirling:

\[
n!\sim\sqrt{2\pi n}\left(\frac ne\right)^n,
\qquad
\sqrt[n]{n!}\sim\frac ne.
\]

Qui \(f\ll g\) significa \(f/g\to0\) nel regime indicato.

Approfondimenti: \href{https://ingegnerismo.it/matematica/limite-di-successione/}{limite di successione} e \href{https://ingegnerismo.it/matematica/numero-di-nepero/}{numero di Nepero}.

\section{Medie di Cesàro e teorema di Stolz--Cesàro (estensione)}\label{medie-di-cesuxe0ro-e-teorema-di-stolzcesuxe0ro-estensione}

\[
a_n\to L
\Longrightarrow
\frac1n\sum_{k=1}^n a_k\to L,
\]

con \(L\in\overline{\mathbb R}\); il converso non vale.

Se \((b_n)\) è strettamente crescente, \(b_n\to+\infty\) e

\[
\lim_{n\to\infty}\frac{a_{n+1}-a_n}{b_{n+1}-b_n}=L
\in\overline{\mathbb R},
\]

allora

\[
\lim_{n\to\infty}\frac{a_n}{b_n}=L.
\]

Approfondimento: \href{https://ingegnerismo.it/matematica/teorema-di-stolz-cesaro/}{teorema di Stolz--Cesàro e medie di Cesàro}.
